\section{Riesz energy of three points on a circle}\label{sec:riesz}

\subsection{The question}
Let $K\subset\R^2$ be compact and $r>0$. Write
\[
U_r(K)\;=\;\max_{\mu}\iint |x-y|^{r}\,\dd\mu(x)\,\dd\mu(y),
\]
the maximum over Borel probability measures $\mu$ on $K$ (a maximiser is called an \emph{equilibrium measure}). In the notation of Fan~\cite{Fan}, $U_r(K)=V_{-r}(K)$ and the Riesz capacity is $\mathrm{Cap}_{-r}(K)=U_r(K)^{1/r}$, so maximising capacity and maximising $U_r$ are the same problem. Fan studies the vertex set $P_N$ of the regular $N$-gon inscribed in the unit circle and states two conjectures, which we quote.

\begin{conjecture}[Fan~\cite{Fan}, Conjecture 5.9]\label{conj:fan59}
Suppose $r\ge2$. Let $A,B,C$ be three points on the unit circle. Fix $B$ and $C$ such that they are symmetric with respect to the $x$-axis, having polar angles $\varphi$ and $-\varphi$, respectively, where $\varphi\in(\pi/2,2\pi/3]$. Move $A$ along the circle with polar angle $\psi\in[0,2\pi-3\varphi]$, so that $BC$ remains the longest side of the triangle $\triangle ABC$. Then the capacity of the three-point set $\{A,B,C\}$ is maximized when $\psi=2\pi-3\varphi$, i.e., when $AC=BC$.
\end{conjecture}

\begin{conjecture}[Fan~\cite{Fan}, Conjecture 5.6]\label{conj:fan56}
Let $r>2$ and let $P_N$ be the vertex set of a regular $N$-gon inscribed in the unit circle ($N\ge3$) with $N$ odd. Let $\theta_N=\frac{(N-1)\pi}{2N}$. Then the energy equals $U_r(P_N)=\dfrac{2^{r-1}\sin^r\theta_N}{1-2^{r-2}\cos^r\theta_N}$, and an equilibrium measure is supported on the three points $1$, $e^{2\theta_N i}$, $e^{2(\pi-\theta_N)i}$ (up to rotation), with masses $\dfrac{2b^r-a^r}{4b^r-a^r},\ \dfrac{b^r}{4b^r-a^r},\ \dfrac{b^r}{4b^r-a^r}$ respectively, where $a=2\sin(\pi/N)$ and $b=2\sin\theta_N$.
\end{conjecture}

Fan proves the special case of Conjecture~\ref{conj:fan59} in which $A$ is replaced by the symmetric configuration (her Lemma 5.8, restated below), shows that Conjecture~\ref{conj:fan59} together with that lemma implies Conjecture~\ref{conj:fan56}, and describes her attempt to prove Conjecture~\ref{conj:fan59} as unsuccessful. We prove both conjectures. The proof is a calculus argument; its only external inputs are three results quoted in~\cite{Fan}, which we use as stated there.

\begin{theorem}\label{thm:riesz}
Conjecture~\ref{conj:fan59} and Conjecture~\ref{conj:fan56} are true. More precisely, in the setting of Conjecture~\ref{conj:fan59}, the energy $U_r(\{A,B,C\})$ is a nondecreasing function of $\psi\in[0,2\pi-3\varphi]$.
\end{theorem}

\noindent\emph{Notation.} In Conjecture~\ref{conj:fan56} the symbols $a,b$ are Fan's; in the proof below they correspond to $\beta$ and $\ell$. Elsewhere $a,b,c$ denote side lengths of a triangle.\medskip

\subsection{Inputs taken from the literature}
\begin{itemize}[leftmargin=2em]
\item[\textbf{(I1)}] \emph{Energy of three points} (Clark--Laugesen~\cite{CL}, Theorem 8, quoted as Lemma 5.7 in~\cite{Fan}). Let $T\subset\R^2$ have three points with mutual distances $a,b,c$, $0<a,b\le c$, and let $r>0$. If $a^r+b^r\le c^r$ then $U_r(T)=c^r/2$; otherwise
\[
U_r(T)=\frac{2a^rb^rc^r}{4a^rb^r-(a^r+b^r-c^r)^2}.
\]
\item[\textbf{(I2)}] \emph{Support} (Bj\"orck~\cite{Bj}, Theorem 12, quoted in~\cite{Fan} before Lemma 5.2). If $r>2$, every compact $K\subset\R^2$ has an equilibrium measure supported on at most three points, which are extreme points of the convex hull of $K$. In particular $U_r(P_N)=\max\{U_r(T): T\subset P_N,\ |T|\le3\}$.
\item[\textbf{(I3)}] \emph{Isosceles triples} (Fan, Lemma 5.8). Let $r\ge2$, $A=(1,0)$, $B=e^{i\alpha}$, $C=e^{-i\alpha}$. Then $U_r(\{A,B,C\})$ is nondecreasing in $\alpha\in[2\pi/3,\pi)$.
\end{itemize}
We did not open the papers of Bj\"orck and Clark--Laugesen; (I1) and (I2) are used exactly as they are stated in~\cite{Fan}.

\subsection{Notation and a formula for the derivative}
Throughout, $r\ge 2$. For a triple with sides $a,b\le c$ put $x=a^r$, $y=b^r$, $z=c^r$ and $s=x+y-z$. If $s\le0$ we are on the \emph{two-point branch}, $U=z/2$; if $s>0$ we are on the \emph{three-point branch}, and then, since $z\ge\max(x,y)$ gives $s\le\min(x,y)$,
\[
D:=4xy-s^2\ \ge\ 4xy-\min(x,y)^2\ >0,\qquad U=\frac{2xyz}{D},\qquad U-\frac z2=\frac{z\,s^2}{2D}>0 .
\]
At $s=0$ both expressions equal $z/2$, so $U$ is a continuous function of $(x,y,z)$.

Fix $\varphi\in(\pi/2,2\pi/3]$ and let $B=e^{i\varphi}$, $C=e^{-i\varphi}$, $A=e^{i\psi}$ with $0\le\psi\le 2\pi-3\varphi$. Put
\[
\lambda=\frac{\varphi-\psi}{2},\qquad \eta=\frac{\varphi+\psi}{2},\qquad\text{so}\qquad 0\le\lambda\le\eta\le\pi-\varphi<\frac\pi2<\varphi .
\]
(The inequality $\eta\le\pi-\varphi$ is $\psi\le2\pi-3\varphi$; $\lambda>0$ for $\psi<\varphi$, which holds because $2\pi-3\varphi<\varphi$ when $\varphi>\pi/2$.) The chord lengths are
\[
a:=|AB|=2\sin\lambda,\qquad b:=|AC|=2\sin\eta,\qquad c:=|BC|=2\sin\varphi .
\]
Since $\sin$ is increasing on $[0,\pi/2]$ and $\sin\varphi=\sin(\pi-\varphi)$, we have $a\le b\le c$, with $b=c$ exactly at $\psi=2\pi-3\varphi$ and $a=b$ exactly at $\psi=0$. So $BC$ is the longest side, as in Conjecture~\ref{conj:fan59}. As $\psi$ increases, $a$ decreases and $b$ increases, $c$ is constant. Write $F(\psi)=U_r(\{A,B,C\})$; $F$ is continuous.

\begin{lemma}[derivative]\label{lem:deriv}
On a $\psi$-interval where $s>0$ and $0<\psi<2\pi-3\varphi$ (so $0<\lambda<\mu$ and $b<c$),
\[
\frac{\dd\log F}{\dd\psi}=\frac{s\,r}{2D}\Big[(z+y-x)\cot\eta-(z+x-y)\cot\lambda\Big].
\]
\end{lemma}
\begin{proof}
From $\log U=\log2+\log x+\log y+\log z-\log D$ and $D=4xy-s^2$, $s=x+y-z$,
\[
\frac{\partial\log U}{\partial x}=\frac1x-\frac{4y-2s}{D}=\frac{D-x(4y-2s)}{xD}=\frac{s(2x-s)}{xD}=\frac{s\,(z+x-y)}{xD},
\]
and symmetrically $\partial_y\log U=s(z+y-x)/(yD)$. Now $\dd\lambda/\dd\psi=-\tfrac12$, $\dd\eta/\dd\psi=\tfrac12$, so $\dd\log x/\dd\psi=-\tfrac r2\cot\lambda$ and $\dd\log y/\dd\psi=\tfrac r2\cot\eta$, while $z$ does not depend on $\psi$. The chain rule gives the formula.
\end{proof}

\subsection{Monotonicity of $F$}
\begin{lemma}[an elementary inequality]\label{lem:ineq}
Let $0<p<q<1$, let $r\ge2$ and suppose $p^r+q^r>1$. Then $q^r-p^r\ \ge\ q^2-p^2$.
\end{lemma}
\begin{proof}
Put $G(\rho)=q^\rho-p^\rho-(q^2-p^2)$ for real $\rho\ge2$, so $G(2)=0$ and $G'(\rho)=q^\rho\ln q-p^\rho\ln p$. Since $\ln p,\ln q<0$, $G'(\rho)\ge0$ is equivalent to $(q/p)^\rho\le \ln p/\ln q$. The left side is strictly increasing in $\rho$ (because $q>p$), so the set $\{\rho\ge2:G'(\rho)\ge0\}$ is an initial segment of $[2,\infty)$ (possibly empty): $G$ is nondecreasing on it and strictly decreasing afterwards. In either case, on every interval $[2,\rho_1]$ we have $G\ge\min\{G(2),G(\rho_1)\}$.

The function $\rho\mapsto p^\rho+q^\rho$ is strictly decreasing and tends to $0$ as $\rho\to\infty$; let $\rho_1$ be the root of $p^{\rho_1}+q^{\rho_1}=1$. Because $p^r+q^r>1$ we have $2\le r<\rho_1$, so it suffices to prove $G(\rho_1)\ge0$.

Put $w=q^{\rho_1}$. Then $p^{\rho_1}=1-w$, $w\in(\tfrac12,1)$ (as $q>p$), and with $\sigma=2/\rho_1\in(0,1]$ we have $q^2=w^\sigma$, $p^2=(1-w)^\sigma$. Thus
\[
G(\rho_1)=(2w-1)-\big(w^\sigma-(1-w)^\sigma\big)=\chi(1)-\chi(\sigma),\qquad \chi(\sigma):=w^\sigma-(1-w)^\sigma ,
\]
and it suffices to show that $\chi$ is nondecreasing on $(0,1]$. We have $\chi'(\sigma)=(1-w)^\sigma|\ln(1-w)|-w^\sigma|\ln w|$, so
\[
\frac{\chi'(\sigma)}{w^\sigma|\ln w|}=\Big(\frac{1-w}{w}\Big)^{\sigma}\frac{|\ln(1-w)|}{|\ln w|}-1\ \ge\ \frac{(1-w)|\ln(1-w)|}{w|\ln w|}-1=\frac{k(1-w)}{k(w)}-1,
\]
where $k(u)=-u\ln u$ and we used $\frac{1-w}{w}<1$, $\sigma\le1$. It remains to see $k(1-w)\ge k(w)$ for $w\in[\tfrac12,1)$. Put $u=1-w\in(0,\tfrac12]$ and $E(u)=k(u)-k(1-u)$. Then $E$ is continuous on $[0,\tfrac12]$ with $E(0)=E(\tfrac12)=0$, and $E''(u)=-\frac1u+\frac1{1-u}=\frac{2u-1}{u(1-u)}<0$ on $(0,\tfrac12)$, so $E$ is concave and $E\ge0$. Hence $\chi'\ge0$.
\end{proof}

\begin{proposition}\label{prop:mono}
$F$ is nondecreasing on $[0,2\pi-3\varphi]$.
\end{proposition}
\begin{proof}
On the open set where $s>0$ and $0<\psi<2\pi-3\varphi$ we use Lemma~\ref{lem:deriv}. Here $s>0$, $D>0$, $\cot\eta>0$ (as $\eta<\pi/2$), $\cot\lambda>0$, and $z+x-y\ge x>0$ (as $y\le z$). So the sign of $\dd\log F/\dd\psi$ is the sign of
\[
R_r-\frac{\cot\lambda}{\cot\eta},\qquad R_r:=\frac{z+y-x}{z+x-y}.
\]
Put $p=a/c$, $q=b/c$ (so $0<p<q<1$) and $t_r=q^r-p^r\in(-1,1)$. Then $R_r=\dfrac{1+t_r}{1-t_r}$, which is an increasing function of $t_r$. For $r=2$ we compute with $\varphi=\lambda+\eta$ and $\sin^2\eta-\sin^2\lambda=\sin(\eta-\lambda)\sin(\eta+\lambda)=\sin\psi\sin\varphi$:
\[
R_2=\frac{\sin^2\varphi+\sin\varphi\sin\psi}{\sin^2\varphi-\sin\varphi\sin\psi}=\frac{\sin\varphi+\sin\psi}{\sin\varphi-\sin\psi}=\frac{2\sin\eta\cos\lambda}{2\cos\eta\sin\lambda}=\frac{\cot\lambda}{\cot\eta},
\]
where $\sin\varphi+\sin\psi=2\sin\eta\cos\lambda$ and $\sin\varphi-\sin\psi=2\cos\eta\sin\lambda$. Hence the sign of $\dd\log F/\dd\psi$ is the sign of $R_r-R_2$, i.e.\ of $t_r-t_2$, i.e.\ of $q^r-p^r-(q^2-p^2)$.

We check the hypotheses of Lemma~\ref{lem:ineq}: $0<p<q<1$ was noted above, and $p^r+q^r>1$ is exactly $s>0$. By Lemma~\ref{lem:ineq}, $t_r\ge t_2$, hence $\dd\log F/\dd\psi\ge0$ wherever $s>0$ in the open interval.

It remains to treat the branches. Suppose $s(\psi_1)>0$ for some $\psi_1\in[0,2\pi-3\varphi)$, and let $\psi_2=\sup\{\psi'\in(\psi_1,2\pi-3\varphi]: s>0\text{ on }[\psi_1,\psi')\}$. On $[\psi_1,\psi_2)$ the function $F$ is continuous and nondecreasing (it is differentiable with nonnegative derivative on the open set where $s>0$ and $0<\psi<2\pi-3\varphi$), so $F(\psi_2)\ge F(\psi_1)>z/2$ by continuity. If $\psi_2<2\pi-3\varphi$ then $s(\psi_2)=0$ and $F(\psi_2)=z/2$, a contradiction; hence $\psi_2=2\pi-3\varphi$, i.e.\ once $s>0$ it stays positive up to the endpoint, and $F$ is nondecreasing on $[\psi_1,2\pi-3\varphi]$. On the set where $s\le0$, $F\equiv z/2$ is constant, and $s>0$ at $\psi$ implies $s>0$ later, so this set is an initial interval $[0,\psi_*]$ (possibly empty) on which $F=z/2$. Since $F$ is continuous and $F\ge z/2$ always, $F$ is nondecreasing on $[0,2\pi-3\varphi]$. The degenerate case $\varphi=2\pi/3$ (where $2\pi-3\varphi=0$) is trivial.
\end{proof}

\begin{proof}[Proof of Conjecture~\ref{conj:fan59}]
By Proposition~\ref{prop:mono}, $F(\psi)\le F(2\pi-3\varphi)$, i.e.\ the energy, and hence the capacity $U_r^{1/r}$, is maximal at $\psi=2\pi-3\varphi$, where $b=c$ (that is $AC=BC$).
\end{proof}

\begin{remark}
The inequality of Lemma~\ref{lem:ineq} gives only $\ge$, so we claim monotonicity (not strict monotonicity) of $F$, which is all that Conjecture~\ref{conj:fan59} asserts. At $r=2$, $t_2-t_2=0$ and $F$ is constant, as in Theorem 5.4 of~\cite{Fan}.
\end{remark}

\subsection{Proof of Conjecture~\ref{conj:fan56}}
Let $N\ge3$ be odd, $r>2$, $\theta=\theta_N=\frac\pi2-\frac\pi{2N}$, and let $v_k=e^{2\pi ik/N}$. Put $\ell=2\sin\theta$ (the longest chord of $P_N$) and $\beta=2\sin(\pi/N)$. Since $\pi/N=\pi-2\theta$ we have $\beta=2\sin2\theta=4\sin\theta\cos\theta$.

\emph{The conjectured triple.} Let $T_N=\{1,\,e^{i(\pi-\pi/N)},\,e^{i(\pi+\pi/N)}\}$; this is the triple of Conjecture~\ref{conj:fan56} because $2\theta_N=\pi-\pi/N$ and $2(\pi-\theta_N)=\pi+\pi/N$. For odd $N$ these are vertices of $P_N$ ($\frac{N\mp1}{2}$ are integers). $T_N$ is isosceles with two sides $\ell$ and base $\beta\le\ell$ (this is $\sin\frac{\pi}{2N}\le\frac12$, true for $N\ge3$). Applying (I1) with $c=\ell$ and the two other sides $\ell$, $\beta$ (so $\ell^r+\beta^r>\ell^r$: three-point branch),
\[
U_r(T_N)=\frac{2\ell^{2r}\beta^r}{4\ell^r\beta^r-\beta^{2r}}=\frac{2\ell^{2r}}{4\ell^r-\beta^r}=\frac{2^{r-1}\sin^r\theta}{1-2^{r-2}\cos^r\theta},
\]
where the last equality uses $\ell=2\sin\theta$, $\beta=4\sin\theta\cos\theta$. This is the energy claimed in Conjecture~\ref{conj:fan56}. Moreover let $\mu_N$ be the measure on $T_N$ with the masses of Conjecture~\ref{conj:fan56} (the apex $1$ carries $\frac{2\ell^r-\beta^r}{4\ell^r-\beta^r}$, and the other two points carry $\frac{\ell^r}{4\ell^r-\beta^r}$ each; here $a=\beta$, $b=\ell$ in the notation of the conjecture). These masses are positive ($2\ell^r>\beta^r$) and sum to $1$, and a direct computation gives
\[
\iint|x-y|^r\dd\mu_N\dd\mu_N=2\Big[2\,m_1m_2\ell^r+m_2^2\beta^r\Big]=\frac{2\ell^{2r}\,(4\ell^r-\beta^r)}{(4\ell^r-\beta^r)^2}=U_r(T_N).
\]

\emph{Upper bound.} By (I2) it suffices to show $U_r(T)\le U_r(T_N)$ for every $T\subset P_N$ with at most three points. Single points give $U_r=0$, and for a pair at distance $d$ one has $U_r=\max_{0\le p\le1}2p(1-p)d^r=d^r/2$, so pairs give $U_r=d^r/2\le\ell^r/2\le U_r(T_N)$ (where $d\le\ell$ is the distance, and $U_r(T_N)\ge\ell^r/2$ because $U\ge z/2$ in (I1)). Let $T=\{A,B,C\}\subset P_N$ be a triple with longest side $BC$.
\begin{itemize}[leftmargin=2em]
\item If $a^r+b^r\le c^r$ (two-point branch), $U_r(T)=c^r/2\le\ell^r/2\le U_r(T_N)$.
\item If $a^r+b^r>c^r$, then, since $(a/c)^r\le(a/c)^2$ and $(b/c)^r\le (b/c)^2$ for $r>2$, we get $a^2+b^2>c^2$, so all angles of the triangle are acute and $T$ is an acute inscribed triangle: its three arcs are all smaller than $\pi$. Rotate so that the midpoint of the arc $BC$ not containing $A$ is $e^{i\pi}$; if that arc has length $2\pi-2\varphi$ (which is $<\pi$, so $\varphi>\pi/2$) then $B,C=e^{\pm i\varphi}$ (after naming), and $A=e^{i\psi}$ lies on the complementary arc, $|\psi|<\varphi$; after a reflection in the $x$-axis, $\psi\in[0,\varphi)$; the arcs $AB,AC$ have lengths $\varphi-\psi,\ \varphi+\psi$; $BC$ longest means $2\pi-2\varphi\ge\varphi+\psi$, i.e.\ $\psi\le2\pi-3\varphi$; and since the three arcs sum to $2\pi$ the longest is at least $2\pi/3$, so $\varphi\le2\pi/3$. Thus $T$ is exactly in the configuration of Theorem~\ref{thm:riesz} and, by Proposition~\ref{prop:mono}, $U_r(T)\le U_r(T')$ where $T'$ is the configuration with $\psi=2\pi-3\varphi$, i.e.\ $|AC|=|BC|=:L=2\sin\varphi$. An isosceles triangle inscribed in the unit circle with legs $L$ is unique up to congruence (its apex is a point of the circle and the other two vertices are the two points of the circle at distance $L$ from it), hence $T'$ is congruent to $\{(1,0),e^{i\alpha},e^{-i\alpha}\}$ with $2\sin(\alpha/2)=L$, i.e.\ $\alpha=2\pi-2\varphi\in[2\pi/3,\pi)$ (as $\varphi\in(\pi/2,2\pi/3]$). As $L=|BC|$ is a chord of $P_N$, $L\le\ell=2\sin\theta$, so $\alpha/2\le\theta$ (both lie in $(0,\pi/2)$) and $\alpha\le2\theta=\pi-\pi/N$, where $2\theta\ge2\pi/3$ because $N\ge3$. By (I3), $U_r(T')\le U_r(\{(1,0),e^{i(\pi-\pi/N)},e^{-i(\pi-\pi/N)}\})=U_r(T_N)$.
\end{itemize}
Hence $U_r(P_N)\le U_r(T_N)=\iint|x-y|^r\dd\mu_N\dd\mu_N\le U_r(P_N)$, so $U_r(P_N)=U_r(T_N)$ and $\mu_N$ is an equilibrium measure of $P_N$, supported on the three points of Conjecture~\ref{conj:fan56}. This proves Conjecture~\ref{conj:fan56}. \hfill$\square$

\begin{remark}
We prove existence of an equilibrium measure with the stated support and masses; we do not claim uniqueness up to rotation (Fan notes the same limitation of her numerical evidence). When $3\mid N$ the equilateral triangle of vertices is an admissible acute triple covered by the case analysis above.
\end{remark}

\subsection{Checks}
The derivative formula, Lemma~\ref{lem:ineq}, and the assembly of $U_r(T_N)$ were checked with computer algebra and high-precision numerics, and the conclusion of Conjecture~\ref{conj:fan56} was compared with a brute-force maximisation over all vertex triples for odd $N\le31$ and several values of $r>2$ (see the repository, \texttt{results/riesz\_triangle}). These checks are not part of the proof.
